% 子文件：4.毛泽东思想及其历史地位，仅input导入
\noparindent
\section{毛泽东思想及其历史地位}

\subsection{毛泽东思想的形成和发展}
\subsubsection{毛泽东思想形成发展的历史条件}
\begin{enumerate}
    \item \textbf{时代背景}：19世纪末20世纪初，世界进入帝国主义和无产阶级革命时代，战争与革命成为时代主题。
    \item \textbf{实践基础}：中国共产党领导人民进行的革命和建设的实践。
\end{enumerate}

\subsubsection{毛泽东思想形成发展的过程}
\paragraph{教材梳理版}
\textbf{时代主题：战争与革命}
\begin{enumerate}
    \item \textbf{开始萌芽：大革命时期}
    \begin{itemize}
        \item 标志性著作：《中国社会各阶级的分析》《湖南农民运动考察报告》等——分析社会各阶级在革命中的地位和作用。
        \item 核心思想：明确无产阶级领导农民是革命核心问题，新民主主义革命思想萌芽。
    \end{itemize}

    \item \textbf{初步形成：土地革命战争时期}
    \begin{itemize}
        \item 标志性著作：《中国的红色政权为什么能够存在?》《井冈山的斗争》《星星之火，可以燎原》《反对本本主义》。
        \item 核心思想：提出农村包围城市、武装夺取政权道路，标志理论初步形成。
    \end{itemize}
    \vspace{0.4em}
    \noindent\sloppy\colorbox{orange!10}{
        \begin{minipage}{0.92\textwidth}
            \textbf{点拨：}中国革命道路：农村建立根据地积蓄力量，时机成熟夺取全国政权。
        \end{minipage}
    }
    \vspace{0.4em}

    \item \textbf{\color{red}成熟：遵义会议后至全民族抗战时期}
    \begin{itemize}
        \item 理论著作：《实践论》《矛盾论》《中国革命和中国共产党》《新民主主义论》等；
        \item 关键节点：1945党的七大将毛泽东思想写入党章，确立指导地位。
    \end{itemize}

    \item \textbf{\color{red}继续发展：解放战争与建国后}
    \begin{itemize}
        \item 标志性著作：七届二中全会报告、《论人民民主专政》《论十大关系》《关于正确处理人民内部矛盾的问题》；
        \item 核心任务：提出马列主义与中国实际“第二次结合”，探索社会主义建设道路。
    \end{itemize}
\end{enumerate}

\paragraph{两次结合拓展}
\begin{itemize}
    \item \textbf{第一次结合}：新民主主义革命阶段，走出革命道路；
    \item \textbf{第二次结合}：社会主义建设阶段，探索建设道路。
\end{itemize}
\vspace{0.3em}
\textbf{“第二次结合”详细阐述}
建国后党完成社会主义改造，进入大规模建设阶段。毛泽东提出把马列主义同中国实际第二次结合，独立探索适合国情的社会主义建设道路，形成人民民主专政、两类矛盾、工业化等理论成果。

\subsection{毛泽东思想的主要内容和活的灵魂}
\subsubsection{毛泽东思想的主题和主要内容}
\begin{enumerate}
    \item \textbf{主题}：中国革命和建设。
    \item \textbf{六大核心内容}：
    \begin{enumerate}
        \item 新民主主义革命理论
        \item 社会主义革命和建设理论
        \item 革命军队与军事战略理论
        \item 政策和策略理论
        \item 思想政治与文化工作理论
        \item 党的建设理论
    \end{enumerate}
\end{enumerate}
\vspace{0.3em}
补充：\keytext{统一战线、武装斗争、党的建设}是革命三大法宝。

\paragraph{拓展：窑洞对}
\textbf{跳出历史周期率两条路径}
\begin{enumerate}
    \item 人民监督政府
    \item 党的自我革命
\end{enumerate}

\subsubsection{毛泽东思想活的灵魂}
\begin{enumerate}
    \item \textbf{总述}：实事求是、群众路线、独立自主；实事求是为精髓。
\end{enumerate}
\vspace{0.4em}
\noindent\sloppy\colorbox{orange!10}{
    \begin{minipage}{0.92\textwidth}
        \textbf{点拨：}1981年历史决议首次概括毛泽东思想活的灵魂。
    \end{minipage}
}
\vspace{0.4em}

\begin{enumerate}[resume]
    \item \textbf{\color{red}实事求是}
    \begin{itemize}
        \item 定位：思想路线、理论精髓；
        \item 内涵：一切从实际出发，理论联系实际，实践检验真理；
        \item 出处：《实践论》《矛盾论》。
    \end{itemize}

    \item \textbf{\color{red}群众路线}
    \begin{itemize}
        \item 定位：党的生命线、根本工作路线；
        \item 内涵：一切为群众、一切依靠群众，从群众中来、到群众中去。
    \end{itemize}

    \item \textbf{\color{red}独立自主}
    \begin{itemize}
        \item 定位：立党立国根本原则；
        \item 出处：1930《反对本本主义》；
        \item 内涵：立足本国力量，自主发展，合理借鉴外部经验。
    \end{itemize}
\end{enumerate}
\vspace{0.4em}
\noindent\sloppy\colorbox{orange!10}{
    \begin{minipage}{0.92\textwidth}
        \textbf{点拨：}《反对本本主义》：中国革命斗争的胜利要靠中国同志了解中国情况。
    \end{minipage}
}
\vspace{0.4em}

\subsection{毛泽东思想的历史地位}
\subsubsection{定义}
毛泽东思想是马列主义在中国的运用发展，是中国革命建设正确理论总结，是党的集体智慧结晶。

\subsubsection{历史地位}
\begin{enumerate}
    \item 马克思主义中国化第一次历史性飞跃理论成果；
    \item 中国革命和建设的科学指南；
    \item 党和人民永久宝贵精神财富。
\end{enumerate}

\subsubsection{正确认识毛泽东与毛泽东思想}
\begin{enumerate}
    \item 评价历史人物必须放在所处时代历史条件下；
    \item 领袖是人不是神，既不能神化崇拜，也不能全盘否定陷入历史虚无；
    \item 总体评价：功绩第一位，错误第二位。
\end{enumerate}

\subsection{本章小结}
毛泽东思想是马克思主义中国化首个重大理论成果，形成于土地革命、成熟于抗战、建国后持续发展。六大理论体系为核心，实事求是、群众路线、独立自主是活的灵魂，长期作为党的指导思想。